\chapter{Introduction}
\label{cap:introducao}

\section{Motivation}

The Linux kernel is a complete operating system, with millions of lines of
code and decades of evolution. Studying how it works from the real code
requires traversing layers of abstraction, compilation configurations,
subsystems and code paths that are not essential to understanding the core of
the problem ``how does a computer go from boot to opening a terminal?''.

The \textbf{Rust VGA Kernel} was created as a didactic slice of that problem.
It reimplements, in Rust, the minimum path necessary to:

\begin{enumerate}
  \item be loaded by a standard bootloader (GRUB via Multiboot2);
  \item leave 32-bit mode and enter x86\_64 long mode;
  \item set up pages and a kernel stack;
  \item initialize an 80$\times$25 VGA text console;
  \item read keys from the PS/2 8042 controller and from the COM1 serial port;
  \item install an IDT, receive exceptions and the PIT timer;
  \item allocate 4 KiB frames and materialize pages on demand;
  \item execute code in ring 3 with user/kernel separation;
  \item offer a small shell with basic commands.
\end{enumerate}

The result is not a reduced Linux, nor a line-by-line port. It is an
independent kernel, with no external Rust dependencies, that reproduces
\emph{concepts} of the equivalent Linux subsystems at a comprehensible scale.

\section{Purpose of the document}

This documentation has four objectives:

\begin{itemize}
  \item \textbf{Explain the architecture} of the kernel: layers, boot flow,
        memory model, execution model and dependencies between modules.
  \item \textbf{Document the main functions} of each module, including
        contracts, preconditions, side effects and limitations.
        Appendix~\ref{ap:funcoes} contains the complete catalog of the
        project's 129 functions.
  \item \textbf{Record the build, execution and validation process}, so that
        another person can reproduce the environment and the tests.
  \item \textbf{Compare the project with the Linux kernel}, presenting
        correspondences of files and functions and a transparent estimate of
        the level of similarity.
\end{itemize}

\section{Scope}

\subsection*{What is included}

\begin{itemize}
  \item BIOS boot from CD/ISO using GRUB and the Multiboot2 protocol.
  \item Transition from 32 bits to long mode (x86\_64).
  \item GDT with ring 0/ring 3 segments, TSS and privilege stacks; static
        page tables at boot and the kernel's higher half.
  \item IDT with 256 gates, exceptions, remapped PIC and PIT timer on IRQ0.
  \item 4 KiB physical frame allocator and \emph{demand paging} for user and
        kernel regions.
  \item Ring 3 demos with U/S protection and a didactic
        \emph{syscall} \code{int \$0x80}.
  \item VGA console in text mode, with hardware cursor, scrolling and color
        attributes.
  \item PS/2 keyboard driver by \emph{polling}, with Shift, Enter, Tab and
        Backspace.
  \item UART 16550 driver on COM1, used for debugging and input.
  \item Shell without dynamic allocation with ten commands.
  \item \emph{Panic handler} that prints the failure and halts the CPU.
  \item Reproducible build with Cargo, \code{ld.lld}, \code{grub-mkrescue}
        and QEMU.
\end{itemize}

\subsection*{What is not included}

\begin{itemize}
  \item Heap, object allocator, \code{alloc} or dynamic collections.
  \item Complete processes, threads, scheduler, context switching, per-process
        address spaces or a broad set of system calls.
  \item Filesystem, networking, USB, ACPI, SMP or native UEFI support.
  \item Drivers for video modes other than the legacy VGA text mode.
  \item Security hardening such as SMEP, SMAP, NX, KASLR or validation of
        user pointers.
\end{itemize}

\section{Target audience}

The document is intended for:

\begin{itemize}
  \item students of operating systems and computer architecture;
  \item Rust developers interested in \code{no\_std} and bare metal;
  \item people who wish to understand the path between BIOS, bootloader and
        console before studying the Linux kernel in depth;
  \item reviewers who need to judge the degree of correspondence between this
        project and Linux.
\end{itemize}

\section{Conventions}

\begin{itemize}
  \item Code names appear in \code{monospaced font}.
  \item Hexadecimal addresses use the \code{0x} prefix.
  \item ``Higher half'' refers to addresses above
        \code{0xFFFF800000000000}, a convention used by Linux and by the
        \code{x86\_64-unknown-none} target.
  \item ``Original'' or ``Linux'' refers, unless stated otherwise, to the Linux
        7.3.1 tree used as the comparison reference.
  \item Address and section size values cited as ``observed'' were extracted
        from the reference compilation and may vary slightly between compiler
        versions.
\end{itemize}

\section{Project status}

The documented version is \code{0.1.0}, with the following characteristics:

\begin{table}[H]
\centering
\caption{Summary of the status of the documented version.}
\label{tab:estado}
\begin{tabularx}{\linewidth}{@{}lX@{}}
\toprule
\textbf{Item} & \textbf{Status} \\
\midrule
Compilation & \texttt{make} completes without errors; target \code{x86\_64-unknown-none}. \\
Image & \code{vga-kernel.iso} El Torito, bootable by BIOS. \\
Boot & Validated in QEMU with SeaBIOS and GRUB Multiboot2. \\
Console & VGA text 80$\times$25, cursor, scrolling and colors. \\
Input & PS/2 by polling and COM1 serial. \\
Interrupts & IDT with 256 gates, exceptions, remapped PIC and PIT at
~100 Hz. \\
Memory & 4 KiB frame allocator, 4 KiB tables and demand paging. \\
Protection & GDT/TSS, ring 3 and U/S bits for user and kernel. \\
Shell & \code{help}, \code{clear}, \code{echo}, \code{info}, \code{demand},
\code{usertest}, \code{userprot}, \code{mem}, \code{reboot}, \code{halt}. \\
Tests & Boot, commands, PS/2/serial input, timer, page faults and ring 3
verified. \\
Linux coverage & Small subset; low similarity (Chapter~\ref{cap:similaridade}). \\
\bottomrule
\end{tabularx}
\end{table}
